\chapter{Thất Bại Học Đường}
\markboth{Thất Bại Học Đường}{Academic Failure}

\section{Rớt Một Kỳ Thi Không Có Nghĩa Cuộc Đời Em Rớt Theo}
\begin{truyen}{Rớt Một Kỳ Thi Không Có Nghĩa Cuộc Đời Em Rớt Theo}{Failing One Exam Does Not Mean Failing Your Life}
\chuhoa{K}{hi trượt nguyện vọng đầu, Quân đóng cửa phòng suốt hai ngày.}
\textit{(When he missed his first-choice admission, Quan shut himself in his room for two days.)}

Em nghĩ mọi thứ kết thúc rồi.
\textit{(He thought everything was over.)}

Bao nhiêu năm học như đổ xuống sông.
\textit{(Years of study seemed thrown away.)}

Người lớn quanh em cũng nói kiểu như tận thế: ``Thế là hỏng.'' Chính những câu đó làm nỗi đau của học sinh phình to hơn sự thật.
\textit{(The adults around him spoke as if it were the apocalypse: ``It is ruined.'' Those sentences made the student’s pain swell larger than reality.)}

Một cậu anh họ kéo Quân đi uống nước và nói thẳng: "Em vừa rớt một cánh cửa, không phải rớt cả cuộc đời.
\textit{(An older cousin took Quan out for a drink and said plainly: "You missed one door, not your entire life.)}

Đừng thần thánh hóa một kỳ thi đến mức để nó quyết định giá trị của em."
\textit{(Do not worship a single exam so much that it gets to define your worth.")}

\end{truyen}

\begin{baihoc}
Thi cử quan trọng, nhưng không đủ quyền để định nghĩa một con người. Sau thất bại, việc đầu tiên là kéo sự việc về đúng kích cỡ thật của nó.
\textit{(Exams matter, but they do not have the authority to define a human being. After failure, the first job is to resize the event back to its true scale.)}
\end{baihoc}

\begin{ghinhoanh}
\textbf{Short English to remember:}

Exams matter, but they do not have the authority to define a human being.

After failure, the first job is to resize the event back to its true scale.

\end{ghinhoanh}

\begin{tuhocnhanh}
\textbf{Cau hoi tu hoc / Self-study questions}

1. Van de chinh la gi? \textit{What is the main problem?}

2. Cach xu ly tot hon la gi? \textit{What is a better action?}

3. Mot cau tieng Anh em muon nho la cau nao? \textit{Which English sentence do you want to remember?}
\end{tuhocnhanh}
\ngancach

\section{Xấu Hổ Là Một Nỗi Đau Thật, Nhưng Không Phải La Bàn}
\begin{truyen}{Xấu Hổ Là Một Nỗi Đau Thật, Nhưng Không Phải La Bàn}{Shame Is Real Pain, But Not a Compass}
\chuhoa{N}{hiều học sinh không sợ điểm thấp bằng sợ bị nhìn thấy điểm thấp.}
\textit{(Many students fear being seen with a low score more than the score itself.)}

Họ sợ ánh mắt bạn bè, giọng họ hàng, câu hỏi hàng xóm, và cả sự im lặng thất vọng trong bữa cơm.
\textit{(They fear classmates’ eyes, relatives’ voices, neighbors’ questions, and the disappointed silence at dinner.)}

Xấu hổ làm người ta muốn trốn.
\textit{(Shame makes people want to hide.)}

Nhưng trốn thì không sửa được gì.
\textit{(But hiding repairs nothing.)}

La bàn của em phải là sự thật và kế hoạch, không phải ánh mắt thiên hạ.
\textit{(Your compass must be truth and planning, not the eyes of the crowd.)}

Một học sinh can đảm nhất không phải người không biết xấu hổ.
\textit{(The bravest student is not the one who feels no shame.)}

Đó là người dám ngồi lại với sự xấu hổ và vẫn tiếp tục làm việc cần làm.
\textit{(It is the one who can sit with shame and still do the work that must be done.)}

\end{truyen}

\begin{baihoc}
Đừng để xấu hổ lái đời em. Cảm xúc đó cần được thừa nhận, nhưng không được giao tay lái.
\textit{(Do not let shame drive your life. That emotion should be acknowledged, but never handed the steering wheel.)}
\end{baihoc}

\begin{ghinhoanh}
\textbf{Short English to remember:}

Do not let shame drive your life.

That emotion should be acknowledged, but never handed the steering wheel.

\end{ghinhoanh}

\begin{tuhocnhanh}
\textbf{Cau hoi tu hoc / Self-study questions}

1. Van de chinh la gi? \textit{What is the main problem?}

2. Cach xu ly tot hon la gi? \textit{What is a better action?}

3. Mot cau tieng Anh em muon nho la cau nao? \textit{Which English sentence do you want to remember?}
\end{tuhocnhanh}
\ngancach

\section{Sau Vấp Ngã, Việc Đầu Tiên Không Phải Là Cố Tỏ Ra Mạnh}
\begin{truyen}{Sau Vấp Ngã, Việc Đầu Tiên Không Phải Là Cố Tỏ Ra Mạnh}{After a Fall, the First Task Is Not to Pretend to Be Strong}
\chuhoa{C}{ó những học sinh vừa thi trượt xong đã đăng trạng thái kiểu: ``Không sao cả.''}
\textit{(Some students fail an exam and instantly post: ``It is fine.'')}

Có thể họ thật sự ổn.
\textit{(Maybe they truly are fine.)}

Nhưng nhiều trường hợp chỉ là đang cố diễn mạnh mẽ để khỏi bị thương hại.
\textit{(But often it is only a performance of strength to avoid pity.)}

Sự thật là có những thất bại đau.
\textit{(The truth is that some failures hurt.)}

Và đau thì phải thừa nhận là đau.
\textit{(And pain should be admitted as pain.)}

Bỏ qua bước đó, em chỉ đang chôn cảm xúc xuống đất chứ không vượt qua nó.
\textit{(If you skip that step, you are burying emotion, not overcoming it.)}

Khóc một trận, nghỉ một ngày, nói chuyện với người tin được, rồi quay lại bàn học.
\textit{(Cry once, rest a day, speak to someone trustworthy, then return to the desk.)}

Đó không phải yếu đuối.
\textit{(That is not weakness.)}

Đó là phục hồi đúng cách.
\textit{(That is recovery done properly.)}

\end{truyen}

\begin{baihoc}
Muốn đứng dậy thật, trước hết phải cho phép mình bị đau thật. Chỉ những thứ được nhìn nhận mới được chữa lành.
\textit{(To stand up for real, you must first permit yourself to hurt for real. Only what is acknowledged can be healed.)}
\end{baihoc}

\begin{ghinhoanh}
\textbf{Short English to remember:}

To stand up for real, you must first permit yourself to hurt for real.

Only what is acknowledged can be healed.

\end{ghinhoanh}

\begin{tuhocnhanh}
\textbf{Cau hoi tu hoc / Self-study questions}

1. Van de chinh la gi? \textit{What is the main problem?}

2. Cach xu ly tot hon la gi? \textit{What is a better action?}

3. Mot cau tieng Anh em muon nho la cau nao? \textit{Which English sentence do you want to remember?}
\end{tuhocnhanh}
\ngancach

\section{Người Bản Lĩnh Không Phải Người Không Thất Bại}
\begin{truyen}{Người Bản Lĩnh Không Phải Người Không Thất Bại}{A Strong Person Is Not One Who Never Fails}
\chuhoa{T}{rường học thường khen người thắng.}
\textit{(Schools often praise winners.)}

Nhưng đời thật cuối cùng lại tôn trọng người biết trở lại sau khi thua.
\textit{(Real life, however, ultimately respects those who know how to return after losing.)}

Một học sinh từng trượt đại học năm đầu, đi làm thêm, học lại, rồi năm sau đỗ vào trường phù hợp hơn.
\textit{(One student failed university admission the first year, worked part-time, studied again, and the next year entered a school that suited him better.)}

Nếu chỉ nhìn một lát cắt, ai cũng bảo em đó thất bại.
\textit{(If you look at one slice of time, everyone calls that student a failure.)}

Nhưng nếu nhìn dài hơn, người ta sẽ thấy em đang học một kỹ năng đắt giá hơn điểm số: khả năng hồi phục.
\textit{(If you look longer, you see a far more valuable skill than a score: recovery.)}

Đời người dài.
\textit{(Life is long.)}

Một cú ngã đau không đáng sợ bằng thói quen ngã xong nằm luôn.
\textit{(One painful fall is less dangerous than the habit of staying down forever.)}

\end{truyen}

\begin{baihoc}
Hãy dạy học sinh không chỉ cách thắng, mà còn cách thua cho đàng hoàng và đứng dậy cho tử tế. Đó mới là bản lĩnh đi đường dài.
\textit{(Teach students not only how to win, but how to lose with dignity and rise with discipline. That is long-distance strength.)}
\end{baihoc}

\begin{ghinhoanh}
\textbf{Short English to remember:}

Teach students not only how to win, but how to lose with dignity and rise with discipline.

That is long-distance strength.

\end{ghinhoanh}

\begin{tuhocnhanh}
\textbf{Cau hoi tu hoc / Self-study questions}

1. Van de chinh la gi? \textit{What is the main problem?}

2. Cach xu ly tot hon la gi? \textit{What is a better action?}

3. Mot cau tieng Anh em muon nho la cau nao? \textit{Which English sentence do you want to remember?}
\end{tuhocnhanh}
\ngancach

\section{Sĩ Tử Trượt Khoa Thi Đầu}
\begin{truyen}{Sĩ Tử Trượt Khoa Thi Đầu}{The Scholar Who Failed His First Exam}
\chuhoa{N}{gày xưa có người trẻ lên kinh ứng thí, chuẩn bị nhiều năm, nhưng trượt ngay kỳ đầu.}
\textit{(Long ago a young scholar traveled to the capital after years of preparation, only to fail at the first examination.)}

Cậu xé hết nháp, định bỏ học về quê làm ruộng.
\textit{(He tore up his drafts and planned to abandon study and return home to farm.)}

Ông chủ quán trọ giữ cậu lại và bảo: "Trượt một khoa không chứng minh con dốt.
\textit{(The innkeeper stopped him and said, "Failing one exam does not prove you are stupid.)}

Nó chỉ chứng minh lần này con chưa đủ."
\textit{(It proves only that this time you were not enough yet.")}

Câu nói đó cứu cậu khỏi quyết định nóng.
\textit{(That sentence saved him from a hot-blooded decision.)}

Ba năm sau cậu thi lại và đỗ.
\textit{(Three years later he retook the exam and passed.)}

\end{truyen}

\begin{baihoc}
Thất bại đầu tiên rất dễ bị ta phóng to thành chân lý cuối cùng. Đó là một ảo giác nguy hiểm.
\textit{(The first failure is easily exaggerated into a final truth. That is a dangerous illusion.)}
\end{baihoc}

\begin{ghinhoanh}
\textbf{Short English to remember:}

The first failure is easily exaggerated into a final truth.

That is a dangerous illusion.

\end{ghinhoanh}

\begin{tuhocnhanh}
\textbf{Cau hoi tu hoc / Self-study questions}

1. Van de chinh la gi? \textit{What is the main problem?}

2. Cach xu ly tot hon la gi? \textit{What is a better action?}

3. Mot cau tieng Anh em muon nho la cau nao? \textit{Which English sentence do you want to remember?}
\end{tuhocnhanh}
\ngancach

\section{Bài Văn Bị Gạch Đỏ Chi Chít}
\begin{truyen}{Bài Văn Bị Gạch Đỏ Chi Chít}{The Essay Covered in Red Marks}
\chuhoa{M}{ột học sinh nhận lại bài Văn với vô số nét đỏ.}
\textit{(A student got back an essay covered in red marks.)}

Em nhìn bài như nhìn bằng chứng cho việc mình vô vọng.
\textit{(She looked at it as if it were proof that she was hopeless.)}

Cô giáo gọi em ở lại và nói: "Cô chữa nhiều vì bài này có thể khá lên.
\textit{(The teacher asked her to stay and said, "I corrected this so much because it can become good.)}

Những bài cô bỏ mặc mới đáng lo."
\textit{(The papers I leave untouched worry me more.")}

Em sững người.
\textit{(She froze.)}

Cùng một tờ giấy, nhưng ý nghĩa đổi hẳn khi được nhìn bằng một lăng kính khác.
\textit{(It was the same paper, but its meaning changed entirely when seen through a different lens.)}

Không phải mọi lời chê đều là phủ định.
\textit{(Not every criticism is a rejection.)}

Có lời chê thực ra là đầu tư.
\textit{(Some criticism is really an investment.)}

\end{truyen}

\begin{baihoc}
Sau thất bại, học sinh cần học cách đọc phản hồi. Không phải chỗ nào đau cũng là chỗ phải bỏ cuộc.
\textit{(After failure, students need to learn how to read feedback. Not every painful place is a place to quit.)}
\end{baihoc}

\begin{ghinhoanh}
\textbf{Short English to remember:}

After failure, students need to learn how to read feedback.

Not every painful place is a place to quit.

\end{ghinhoanh}

\begin{tuhocnhanh}
\textbf{Cau hoi tu hoc / Self-study questions}

1. Van de chinh la gi? \textit{What is the main problem?}

2. Cach xu ly tot hon la gi? \textit{What is a better action?}

3. Mot cau tieng Anh em muon nho la cau nao? \textit{Which English sentence do you want to remember?}
\end{tuhocnhanh}
\ngancach

\section{Năm Gap Không Đáng Xấu Hổ}
\begin{truyen}{Năm Gap Không Đáng Xấu Hổ}{A Gap Year Is Not Shameful}
\chuhoa{M}{ột cậu học sinh thi xong không đỗ và quyết định nghỉ một năm để học lại, làm thêm, và nghĩ kỹ hơn về đường mình đi.}
\textit{(A student failed an entrance exam and decided to take a year to study again, work part-time, and think more carefully about his path.)}

Người quanh cậu nói nhiều câu khiến cậu nhục: chậm một năm, thua bạn bè, tụt hậu.
\textit{(People around him said things meant to humiliate him: one year behind, defeated by friends, falling behind.)}

Nhưng trong năm đó, cậu lớn lên nhiều hơn ba năm ngồi học chỉ để khỏi bị chê.
\textit{(Yet in that year he grew more than in three years of studying only to avoid criticism.)}

Có những đường vòng không phải lạc đường.
\textit{(Some detours are not getting lost.)}

Nó là khoảng cần thiết để người ta đi lại cho đúng.
\textit{(They are necessary intervals for someone to get back on the right road.)}

\end{truyen}

\begin{baihoc}
Không phải lúc nào đi nhanh hơn cũng tốt hơn. Với một số người, chậm lại một năm là cái giá hợp lý để không lạc cả chục năm.
\textit{(Moving faster is not always moving better. For some people, slowing down for a year is a fair price to avoid being lost for ten.)}
\end{baihoc}

\begin{ghinhoanh}
\textbf{Short English to remember:}

Moving faster is not always moving better.

For some people, slowing down for a year is a fair price to avoid being lost for ten.

\end{ghinhoanh}

\begin{tuhocnhanh}
\textbf{Cau hoi tu hoc / Self-study questions}

1. Van de chinh la gi? \textit{What is the main problem?}

2. Cach xu ly tot hon la gi? \textit{What is a better action?}

3. Mot cau tieng Anh em muon nho la cau nao? \textit{Which English sentence do you want to remember?}
\end{tuhocnhanh}
\ngancach

\section{Khi Cả Lớp Biết Em Trượt}
\begin{truyen}{Khi Cả Lớp Biết Em Trượt}{When the Whole Class Knows You Failed}
\chuhoa{Đ}{iều đau nhất đôi khi không phải trượt.}
\textit{(Sometimes the worst pain is not failing itself.)}

Điều đau nhất là sáng hôm sau phải bước vào lớp với cảm giác ai cũng biết.
\textit{(It is walking into class the next morning feeling that everyone knows.)}

Một học sinh đã định nghỉ học mấy hôm chỉ để tránh ánh mắt người khác.
\textit{(One student planned to skip school for several days just to avoid other people’s eyes.)}

Cô giáo nói với em: "Em càng trốn, chuyện này càng lớn.
\textit{(The teacher told him, "The more you hide, the larger this becomes.)}

Em đi học lại đi, để não em thấy rằng em vẫn sống tiếp được."
\textit{(Go back to school, so your mind sees that life continues.")}

Ngày đầu rất khó, nhưng qua được ngày đầu, cái xấu hổ mất bớt quyền lực.
\textit{(The first day was very hard, but once it was crossed, shame lost some of its power.)}

\end{truyen}

\begin{baihoc}
Một phần phục hồi là quay lại nơi mình muốn trốn. Nếu cứ lẩn mãi, nỗi sợ sẽ lớn hơn sự thật.
\textit{(Part of recovery is returning to the place you want to avoid. If you keep hiding, fear grows larger than reality.)}
\end{baihoc}

\begin{ghinhoanh}
\textbf{Short English to remember:}

Part of recovery is returning to the place you want to avoid.

If you keep hiding, fear grows larger than reality.

\end{ghinhoanh}

\begin{tuhocnhanh}
\textbf{Cau hoi tu hoc / Self-study questions}

1. Van de chinh la gi? \textit{What is the main problem?}

2. Cach xu ly tot hon la gi? \textit{What is a better action?}

3. Mot cau tieng Anh em muon nho la cau nao? \textit{Which English sentence do you want to remember?}
\end{tuhocnhanh}
\ngancach

\section{Thua Một Keo, Học Một Nếp}
\begin{truyen}{Thua Một Keo, Học Một Nếp}{Lose One Round, Learn One Discipline}
\chuhoa{M}{ột học sinh trượt đội tuyển và sau đó nhìn lại mới thấy mình thua không chỉ vì kiến thức.}
\textit{(A student failed to enter the elite team and later realized the loss was not about knowledge alone.)}

Em ngủ muộn, làm việc theo hứng, và chỉ tăng tốc sát ngày.
\textit{(He slept late, worked by mood, and accelerated only near deadlines.)}

Thất bại lần này bóc ra thói quen mà lời nhắc nhở trước đó chưa bóc được.
\textit{(This failure exposed habits that previous advice had never truly exposed.)}

Có những cú thua dạy bài học mà lúc còn thắng người ta không chịu học.
\textit{(Some defeats teach lessons that people refuse to learn while winning.)}

Năm sau em chưa chắc thắng ngay, nhưng em bước vào với một nếp sống nghiêm túc hơn hẳn.
\textit{(The next year he did not necessarily win immediately, but he entered with a far more disciplined way of living.)}

\end{truyen}

\begin{baihoc}
Thất bại đáng giá nhất là thất bại lộ ra được thói quen xấu. Nó đau, nhưng nó mở cửa sửa từ gốc.
\textit{(The most valuable failure is the one that exposes bad habits. It hurts, but it opens the door to root-level change.)}
\end{baihoc}

\begin{ghinhoanh}
\textbf{Short English to remember:}

The most valuable failure is the one that exposes bad habits.

It hurts, but it opens the door to root-level change.

\end{ghinhoanh}

\begin{tuhocnhanh}
\textbf{Cau hoi tu hoc / Self-study questions}

1. Van de chinh la gi? \textit{What is the main problem?}

2. Cach xu ly tot hon la gi? \textit{What is a better action?}

3. Mot cau tieng Anh em muon nho la cau nao? \textit{Which English sentence do you want to remember?}
\end{tuhocnhanh}
\ngancach

\section{Không Ai Cười Em Lâu Đến Thế}
\begin{truyen}{Không Ai Cười Em Lâu Đến Thế}{No One Laughs at You for That Long}
\chuhoa{S}{au một lần vấp, nhiều học sinh nghĩ cả thế giới đang nhớ và nhắc đến mình.}
\textit{(After one stumble, many students think the whole world remembers and keeps talking about them.)}

Sự thật phũ hơn nhưng cũng dễ thở hơn: phần lớn người khác bận lo chuyện của chính họ.
\textit{(The truth is harsher but also easier to breathe with: most people are busy with their own lives.)}

Một cậu từng bị xấu hổ vì thuyết trình hỏng.
\textit{(A boy once felt humiliated after a bad presentation.)}

Một tuần sau em phát hiện gần như không ai còn nhắc nữa.
\textit{(One week later he found that almost nobody mentioned it anymore.)}

Chính em là người giữ sự cố đó sống lâu nhất trong đầu mình.
\textit{(He himself was the one keeping the incident alive longest in his own mind.)}

\end{truyen}

\begin{baihoc}
Nhiều cú ngã kéo dài không phải vì thiên hạ bám lấy em, mà vì em bám lấy nó. Buông dần đi là một phần của đứng dậy.
\textit{(Many falls last not because the world clings to them, but because you do. Letting go is part of standing up.)}
\end{baihoc}

\begin{ghinhoanh}
\textbf{Short English to remember:}

Many falls last not because the world clings to them, but because you do.

Letting go is part of standing up.

\end{ghinhoanh}

\begin{tuhocnhanh}
\textbf{Cau hoi tu hoc / Self-study questions}

1. Van de chinh la gi? \textit{What is the main problem?}

2. Cach xu ly tot hon la gi? \textit{What is a better action?}

3. Mot cau tieng Anh em muon nho la cau nao? \textit{Which English sentence do you want to remember?}
\end{tuhocnhanh}
